\section{Introduction}

\subsection*{Thrombolysis in acute ischaemic stroke}

Intravenous thrombolysis has been a cornerstone of acute ischaemic stroke treatment since the NINDS rt-PA trial showed that alteplase given within 3 hours of onset increased the likelihood of minimal or no disability at 3 months, at the cost of an increased risk of symptomatic intracranial haemorrhage.\cite{ninds1995} ECASS III extended the treatment window to 4.5 hours, with 52.4\% of alteplase-treated patients achieving a favourable outcome compared with 45.2\% of those given placebo.\cite{hacke2008} IST-3, the largest single trial, randomised 3035 patients within 6 hours of onset, more than half of whom were aged over 80 years; it did not show a significant difference in its primary outcome but did show a favourable shift in the distribution of functional outcomes.\cite{ist32012}

An individual patient data meta-analysis of 6756 patients from nine randomised trials confirmed that alteplase increased the odds of a good outcome (modified Rankin Scale [mRS] score 0--1), with benefit greatest when treatment was given early (odds ratio [OR] 1.75 when given within 3 hours, 1.26 at 3--4.5 hours and 1.15 beyond 4.5 hours).\cite{emberson2014} The proportional benefit was similar irrespective of age or baseline stroke severity, but alteplase increased the risk of parenchymal haemorrhage within 7 days (6.8\% vs 1.3\%) and of fatal intracranial haemorrhage (2.7\% vs 0.4\%).\cite{emberson2014} Further analyses of the same trials showed that alteplase shifted the whole distribution of functional outcomes towards less disability,\cite{lees2016} and a Cochrane review of 27 trials reached similar conclusions.\cite{wardlaw2014} Since then, imaging selection has allowed treatment of patients with unknown time of onset,\cite{thomalla2018} and tenecteplase has been shown to be non-inferior to alteplase.\cite{menon2022} Guidelines in North America and Europe recommend intravenous thrombolysis for eligible patients with disabling ischaemic stroke within 4.5 hours of onset.\cite{powers2019,berge2021}

\subsection*{Thrombolysis in mild stroke}

Mild stroke is common. More than half of patients with acute ischaemic stroke present with a National Institutes of Health Stroke Scale (NIHSS) score of 0--5,\cite{khatri2018} and mild or rapidly improving symptoms have been one of the most frequently cited reasons for withholding thrombolysis.\cite{smith2011} A mild deficit does not, however, guarantee a good outcome. In the US Get With The Guidelines-Stroke registry, 28\% of patients who were not thrombolysed solely because their stroke was mild or improving were not discharged home, and a similar proportion could not walk without assistance at discharge.\cite{smith2011} Among patients with an NIHSS score of 5 or less who were thrombolysed in routine practice, symptomatic intracranial haemorrhage occurred in 1.8\%, but about 30\% could not return directly home or walk independently at discharge.\cite{romano2015} Most studies define mild (or minor) stroke as an NIHSS score of 0--5; the evidence below therefore refers to this range unless otherwise stated, rather than the NIHSS $<$5 threshold used in this study.

Evidence on thrombolysis in mild stroke from the pivotal trials is limited because few patients with clearly non-disabling deficits were enrolled.\cite{khatri2018} The pooled analysis of the major trials included patients with an NIHSS score of 0--4 and did not find that the proportional benefit of alteplase differed by stroke severity.\cite{emberson2014} In a post hoc analysis of IST-3 restricted to 106 patients treated within 3 hours with an NIHSS score of 5 or less, alteplase was associated with a higher proportion of patients alive and independent at 6 months (84\% vs 65\%; adjusted OR 3.31, 95\% confidence interval [CI] 1.24--8.79), though with no significant effect on the stricter outcome of no or minimal symptoms.\cite{khatri2015}

The randomised trials designed specifically for mild stroke have focused largely on patients whose deficits are judged not to be clearly disabling, and have not shown benefit. PRISMS randomised 313 patients with an NIHSS score of 0--5 and non-disabling deficits to alteplase or aspirin within 3 hours. It was stopped early and found no difference in excellent outcome (mRS 0--1) at 90 days (78.2\% vs 81.5\%; adjusted risk difference $-$1.1\%, 95\% CI $-$9.4\% to 7.3\%), with symptomatic intracranial haemorrhage in 3.2\% of alteplase-treated patients and none of the aspirin group.\cite{khatri2018} Patients in PRISMS had milder deficits than mild-stroke patients treated with alteplase in routine practice, and were comparable to those who were not treated.\cite{asdaghi2021} The ARAMIS trial randomised 760 patients with minor non-disabling stroke within 4.5 hours and found dual antiplatelet therapy to be non-inferior to alteplase (excellent outcome 93.8\% vs 91.4\%).\cite{chen2023} TEMPO-2 randomised 886 patients with an NIHSS score of 0--5 and a proven intracranial occlusion or focal perfusion abnormality to tenecteplase or non-thrombolytic standard care within 12 hours. It was stopped for futility: 72\% of the tenecteplase group and 75\% of the control group returned to their baseline level of function (risk ratio 0.96, 95\% CI 0.88--1.04), and more patients died after tenecteplase (5\% vs 1\%).\cite{coutts2024}

Meta-analyses have reached similar conclusions. Analyses of patients with minor non-disabling stroke found no benefit of alteplase over no thrombolysis,\cite{alhazzani2024} and a network meta-analysis ranked dual antiplatelet therapy above alteplase for excellent functional outcome (OR 1.52, 95\% credible interval 1.09--2.35) and symptomatic intracranial haemorrhage.\cite{lun2025} A meta-analysis of 20 randomised and observational studies (13,397 patients with NIHSS 0--5) found no improvement in functional outcome with thrombolysis but a higher risk of symptomatic intracranial haemorrhage (OR 7.48, 95\% CI 3.55--15.76).\cite{zhang2024} The most recent meta-analysis, restricted to four randomised trials enrolling only minor stroke (3364 patients), found that thrombolysis did not improve the odds of excellent recovery (OR 0.85, 95\% CI 0.70--1.03), and was associated with lower odds of functional independence and higher odds of symptomatic intracranial haemorrhage (OR 5.22) and death at 90 days (OR 2.40).\cite{doheim2025}

Current guidelines therefore recommend against thrombolysis for patients with minor non-disabling stroke, while still recommending or suggesting it for those with mild but disabling deficits.\cite{powers2019,berge2021} The evidence for this distinction is weak. Most patients enrolled in the mild-stroke trials had non-disabling deficits, and in a secondary analysis of TEMPO-2 the 100 patients with disabling deficits did not benefit from tenecteplase either (return to baseline function 54.7\% vs 68.1\%; adjusted risk ratio 0.81, 95\% CI 0.60--1.10).\cite{zhang2025} Whether thrombolysis benefits any patients with mild stroke, and if so which ones, therefore remains uncertain.
